\section{The two filling summands}\label{sec:filling-summands}

The integral lattice is now fixed by \eqref{eq:tube-1}. We compute
its ordered filling form on the two rational summands, first by
separating inner and exterior fillings and then by examining the
normal disks of an inner filling. Throughout, the first filling
bounds the first cycle on the reference page, and the second bounds
the positive neighboring-page transport of the second cycle. The boundary continuations
preserve these complex page values and arise from oriented ambient
isotopies, so they preserve this order and the intersection numbers.

\begin{lemma}[orthogonality and exterior form]\label{lem:exterior}
The \(K_\Q\) and \((H\otimes E)_\Q\) summands are
orthogonal for the tube filling form. The
second summand has form \(s_h\otimes e\).
\end{lemma}
\begin{proof}
The inner classes bound in the acyclic
\(\widehat W\). For an exterior page class
representing an element of \(H\otimes\ker\tau\),
its image in \(H_3(\widehat M_{\rm ext})\)
may be nonzero. The compatible sublevel continuations take \(\widehat J\) to
the product of the acyclic base total of \(h\) with \(L\).
Its integral \(H_3\) is therefore generated by one oriented lens
fiber. Lemma~\ref{lem:acyclic} then gives
\(H_3(\widehat J)\cong H_3(\widehat M_{\rm ext})\cong\Z\). Under the interface continuation, the
inclusion \(J\hookrightarrow\widehat J\) sends the oriented
lens-fiber generator of \(H_3(J)\) to a generator of
\(H_3(\widehat J)\). Subtract the
corresponding integral multiple of the
page-interface generator, which is zero
in the full page. The adjusted exterior
class bounds in \(\widehat M_{\rm ext}\).
This can be done compatibly under a short
value-circle transport arc. Put the two
fillings on opposite sides of a bicollar
of \(\widehat J\); their interiors are then
disjoint, proving orthogonality.

Continue the exterior piece to \(s=0\).
Its full lift to \((U,V)\) lies off the
cone vertex, where the cyclic map is an
orientation-preserving unramified \(k\)-cover.
Thus intersections of full lifts are \(k\)
times intersections downstairs. This covering count uses the
original fillings off the cone vertex. In the contractible
upstairs total, the filling pairing is unchanged when those
fillings are replaced by the radial fillings used below.
The smooth join calculation follows Sakamoto's product formula
\cite[Theorem~2, p.~715]{Sakamoto1974} and the product construction of
\cite{KauffmanNeumann1977}; the cone extension
and integral gluing were established here. The lifted cycles
also live in the full smooth tube
for \(h+\widetilde P\), whose total is
contractible by weights. Its page
Mayer--Vietoris coordinate is
\(H\otimes H_1(\widetilde Q)\);
the lift of \(z\in H\otimes E\) has
coordinate \((1\otimes\tr)z\).
For pure tensors, represent this class
by the join of loops in the two factor
pages, varying their values as
\((tc,(1-t)c)\), \(0\le t\le1\).
It meets the equal-value cut in the
product of the two loops. Fill the join
by the two factor radial cone fillings
with values in
\(\{(t,t')c:t,t'\ge0,\ t+t'\le1\}\).
The two ordered fillings meet only at
the common origins of the factors
because their neighboring rays are
disjoint away from zero. Near each
intersection they are products of
two real 2-dimensional filling
intersections. Their Koszul sign is
\((-1)^{2\cdot2}=+1\); any uniform
join-orientation sign occurs in both
fillings and cancels. Hence the upstairs
form is \(s_h\otimes\widetilde s\).
Division by \(k\) gives \(s_h\otimes e\)
downstairs. Since \(kE\subseteq\ker\tau\), exterior classes span their
rational summand, so the identity holds
on all of it.
\end{proof}

The exterior calculation has now fixed both cross terms and the
product form. The remaining calculation compares an inner filling
with the plane filling that supplies its class in \(K\).

\begin{lemma}[the inner negative sign]\label{lem:inner}
The restriction of the filling form to
\(K\) is \(-s_d|_K\).
\end{lemma}
\begin{proof}
Starting from two plane fillings, we build relative chains and sweep
them by the normal circle at two distinct balancing radii.
The sweeps will be separated away from the root
locus, so each remaining intersection is an intersection of the
plane fillings multiplied by the sign of two normal disks.
We first construct the sweeps as filling chains and then compute
that normal sign.

Use the \(h\)-model of Lemma~\ref{lem:continuations} over the value disk
\(\{|v|\le\delta\}\). For \(|v|\le\delta\), set
\[
G_v=(\mathcal B\cap\{h=v\})\times\{|w|\le\lambda\},\qquad
\widehat G=(\mathcal B\cap\{|h|\le\delta\})\times\{|w|\le\lambda\},
\]
and let \(\widehat\D=\widehat G\cap\{R_s=0\}\) be its root total.
By the selected plane cut calculation used in
Lemma~\ref{lem:acyclic}, \(\widehat\D\) is acyclic; the same holds
for \(\widehat G\), the base total times a disk.

Choose the initial chains as finite integral sums of mapped oriented
simplices, piecewise smooth after compatible subdivision. The swept
chains will be continuous at the conical strata described below and
piecewise smooth elsewhere. Intersections mean transverse intersections of
the maps on their smooth common contacts, counted with their integer
coefficients; we retain set notation for their images.
For two classes \(l_1,l_2\in K\), take \(c_1\) as the reference
page and \(c_2\) as its positive neighbor, and represent \(l_i\)
by cycles on the root page in \(G_{c_i}\). Their images vanish in
\(H_1(G_{c_i})\) because \(l_i\in K=\ker\pi_*\).
This vanishing and the acyclicity of \(\widehat\D\) let us choose
surface chains \(\Sigma_i\subset\widehat\D\) with
\(\partial\Sigma_i=l_i\) and \(V_i\subset G_{c_i}\) with
\(\partial V_i=-l_i\). Arrange \(\Sigma_1,\Sigma_2\)
transversely in the real four-dimensional \(\widehat\D\),
with generic \(V_i\) and disjoint page-boundary collars.
Since \(H_2(\widehat G)=0\),
the closed chain \(\Sigma_i+V_i\) bounds a
\(3\)-chain \(\Gamma_i\subset\widehat G\).
Prescribe a collar of the \(\Sigma_i\)-face on which
\(R_s=t\ge0\) is a positive-real normal coordinate.
Near \(l_i\), make its normal lift tangent to \(h=c_i\),
possible because \((h,R_s)\) is a submersion there, and join
it to a collar of \(V_i\). Keep these collars fixed.
Perturb only the remaining interiors of \(\Gamma_1,\Gamma_2\)
relative to their boundary chains, and make
\[
I_i=(\Gamma_i\setminus\text{prescribed collar})\cap\widehat\D
\]
transverse there. Each \(I_i\) is a real one-dimensional
chain, with any endpoints on the prescribed \(V_i\) contacts.
Relative transversality gives
\(I_1\cap\Sigma_2=\Sigma_1\cap I_2=I_1\cap I_2=\varnothing\):
the respective dimension sums are \(1+2<4\) and \(1+1<4\).
Separate their finite boundary contacts in the distinct
page-value collars. Thus the only common root contacts are
the prescribed \(\Sigma_1\cap\Sigma_2\) intersections.

Choose distinct \(\sigma_1,\sigma_2\) sufficiently near
\(\lambda^{(\alpha-\beta)/2}\) that the strict balancing bounds
in Lemma~\ref{lem:acyclic} hold for both; their circles then stay
inside \(Y_\lambda\).
Sweep each \(\Gamma_i\) by the \(y\)-angle
using \(|y|=\sigma_i\sqrt{|R_s|}\),
\(|z|=\sigma_i^{-1}\sqrt{|R_s|}\), and \(z=R_s/y\)
away from the root.
This is the chain-level relative suspension
\(C_3(\widehat G,\widehat\D)\to C_4(\widehat W)\).
To realize it by mapped chains, triangulate the domains compatibly
with the prescribed collars, the transverse root curves, and their
boundary faces. Sweep each oriented top simplex, collapse its
circle over the root, and push the resulting top chain into
\(\widehat W\). The internal faces cancel with the original
chain coefficients. The \(\Sigma_i\)-face is zero in the relative
complex, so it contributes no boundary after collapse.
Write \(\mathfrak S\) also for the circle-first suspension of
Section~\ref{sec:circle-collapse} followed by this collapse and
pushforward. Give the sweep the negative of the raw product
orientation. In the relative quotient chain complex the two
relevant signs are
\[
 \partial\bigl(-\mathfrak S(\Gamma_i)\bigr)
    =+\mathfrak S(V_i),\qquad
 -\partial_{\D}V_i=l_i.
\]
The second identity says that this boundary is the positive
\(K\)-page class in the convention fixed after \eqref{eq:tube-12}.
The same reversal for both fillings does not change their
intersection product.

We now verify that the mapped sweeps have the claimed filling
boundary at root contacts.
Along the prescribed normal collar,
the collapsed circle is the center of a disk. At an incidental
transverse root curve \(I_i\), the local quotient is
\(I_i\times\operatorname{cone}(T^2)\): its torus link is closed
and oriented. The oriented top simplices of a triangulated cone
have cancelling internal faces, and its tip is not a
codimension-one face; hence this local model contributes no
additional chain boundary.
The map extends continuously to \(y=z=0\) there.

We next locate the intersections of the two sweeps.
Where \(R_s\ne0\), the sweeps are
disjoint at a common base point because their \(|y|\) differ.
Where \(R_s=0\), their possible common contacts are
\((\Sigma_1\cup I_1)\cap(\Sigma_2\cup I_2)
=\Sigma_1\cap\Sigma_2\). Thus all intersections occur
above the prescribed transverse surface intersections.
In particular, the conical tips over \(I_i\) are disjoint from
the other swept chain. A small general-position approximation in
disjoint neighborhoods of those tips, fixed on the prescribed
collars and with each boundary kept in its assigned page, preserves
the filling classes and adds no intersections. Thus the filling
intersection is computed entirely on the smooth prescribed collars.
Only the intersection sign of the normal disks remains.
In complex normal coordinates \(y,z\)
there, each swept collar is the real
disk \(z=\sigma_i^{-2}\overline y\).
In real coordinates \(y=x+it\), its
graph map has matrix
\(\mathrm{diag}(\sigma_i^{-2},
-\sigma_i^{-2})\), so the ordered
intersection determinant is
\[
\det\bigl(
\mathrm{diag}(\sigma_2^{-2},-\sigma_2^{-2})
-\mathrm{diag}(\sigma_1^{-2},-\sigma_1^{-2})
\bigr)
=-(\sigma_2^{-2}-\sigma_1^{-2})^2<0.
\]
Both collar disks have the same orientation relative to
the complex \(y\)-orientation (both are reversed by the
chosen sweep), so their common reversals cancel in the
intersection product. Also,
interchanging a disk past a root
surface contributes
\((-1)^{2\cdot2}=+1\).
Thus every root-surface
intersection sign is reversed.
The root-surface filling
intersections are exactly
\(s_d(l_1,l_2)\) by the common
nearby-page convention, proving
the lemma. Changing the
uniform sign of the \(K\)-coordinate
does not change this bilinear
restriction.
\end{proof}

\noindent\emph{Proof of Theorem~\ref{thm:tube}.}
Lemmas~\ref{lem:boundary}, \ref{lem:continuations}, and
\ref{lem:isolation} provide a
proper Milnor-model tube with
a unique critical point. The acyclicity
proved in Lemma~\ref{lem:acyclic} completes
the comparison in Lemma~\ref{lem:local},
identifying its page and pairing with the
ordinary local Milnor data.
Lemmas~\ref{lem:acyclic},
\ref{lem:mv}, \ref{lem:bock},
and \ref{lem:lattice} give the
integral lattice \eqref{eq:tube-1}, while
Lemmas~\ref{lem:exterior} and
\ref{lem:inner} give \eqref{eq:tube-2}.
\qed
